\documentclass{article}
\usepackage{graphicx} % Required for inserting images


% -------------------- Packages --------------------
\usepackage[utf8]{inputenc}
\usepackage[T1]{fontenc}
\usepackage{lmodern}          % Better font
\usepackage{geometry}         % Margins
\geometry{margin=1in}

\usepackage{amsmath, amssymb, amsthm}   % Math packages + theorem env
\usepackage{mathtools}        % Enhancements for amsmath
\usepackage{enumitem}         % Better lists
\usepackage{graphicx}         % Images
\usepackage{booktabs}         % Good tables
\usepackage{hyperref}         % Clickable refs
\usepackage{xcolor}           % Colors, if needed
\usepackage{tikz}             % Graph drawing (optional)
\usetikzlibrary{graphs}       % Graph theory tools (optional)

% -------------------- Theorem Environments --------------------
\theoremstyle{plain}
\newtheorem{theorem}{Theorem}
\newtheorem{lemma}{Lemma}
\newtheorem{proposition}{Proposition}
\newtheorem{corollary}{Corollary}

\theoremstyle{definition}
\newtheorem{definition}{Definition}
\newtheorem{example}{Example}

\theoremstyle{remark}
\newtheorem{remark}{Remark}

% Proof environment already included in amsthm.

% -------------------- Custom Commands --------------------
\newcommand{\N}{\mathbb{N}}
\newcommand{\Z}{\mathbb{Z}}
\newcommand{\Q}{\mathbb{Q}}
\newcommand{\R}{\mathbb{R}}
\newcommand{\C}{\mathbb{C}}

% Handy graph notation
\newcommand{\vn}{V(G)}
\newcommand{\en}{E(G)}

% ---------------------------------------------------------

\title{}
\author{}
\date{}

\begin{document}

\maketitle

\section{First Order Logic}

% Hilbert axioms for implication

\[
\begin{array}{ll}
\text{(A1)} & A \to (B \to A)\\[2mm]
\text{(A2)} & (A \to (B \to C)) \to ((A \to B) \to (A \to C)) \\[2mm]
\text{(A3)} & (\neg B \to \neg A) \to (A \to B)
\end{array}
\]

\noindent\textbf{Claim.} \(\;\vdash P \to P.\)

\begin{proof}
\[
\begin{array}{rll}
(1) & P \to (P \to P) 
    & \text{instance of (A1) with } A := P,\; B := P \\[2mm]

(2) & P \to ((P \to P) \to P) 
    & \text{instance of (A1) with } A := P,\; B := (P \to P) \\[2mm]

(3) & (P \to ((P \to P) \to P)) \to ((P \to (P \to P)) \to (P \to P))
    & \text{instance of (A2) with } A := P,\; B := (P \to P),\; C := P \\[2mm]

(4) & (P \to (P \to P)) \to (P \to P)
    & \text{from (2) and (3) by Modus Ponens} \\[2mm]

(5) & P \to P
    & \text{from (1) and (4) by Modus Ponens}
\end{array}
\]
Hence \(\vdash P \to P\).
\end{proof}

\noindent\textbf{Claim.} \(\;\{\,P \to Q,\; Q \to R\,\} \;\vdash\; P \to R.\)

\begin{proof}
\[
\begin{array}{rll}
(1) & (Q \to R) \to (P \to (Q \to R))
    & \text{instance of (A1) with } A := (Q \to R),\; B := P \\[2mm]

(2) & Q \to R
    & \text{given} \\[2mm]

(3) & P \to (Q \to R)
    & \text{from (1) and (2) by Modus Ponens} \\[2mm]

(4) & (P \to (Q \to R)) 
      \to \big((P \to Q) \to (P \to R)\big)
    & \text{instance of (A2) with } 
      A := P,\; B := Q,\; C := R \\[2mm]

(5) & (P \to Q) \to (P \to R)
    & \text{from (3) and (4) by Modus Ponens} \\[2mm]

(6) & P \to Q
    & \text{given} \\[2mm]

(7) & P \to R
    & \text{from (5) and (6) by Modus Ponens}
\end{array}
\]
Thus \(\;P \to R\) is derivable from the assumptions.
\end{proof}

\noindent\textbf{Claim.} The formulas underneath are logically equivalent
\[
\forall x\,\forall y\,(q(x,y)\to p(x))
\quad\text{ and}\quad
\forall x\,(\exists y\,q(x,y)\to p(x))
\]


\begin{proof}
\[
\begin{aligned}
\forall x\forall y\,(q(x,y)\to p(x))
&\equiv \forall x\forall y\,(\neg q(x,y)\lor p(x))\\
&\equiv \forall x\bigl(\forall y(\neg q(x,y)\lor p(x))\bigr)\\
&\equiv \forall x\bigl(\neg\exists y\,q(x,y)\lor p(x)\bigr)\\
&\equiv \forall x\bigl(p(x)\lor \neg\exists y\,q(x,y)\bigr)\\
&\equiv \forall x\bigl(\exists y\,q(x,y)\to p(x)\bigr).
\end{aligned}
\]
Thus the two statements are equivalent.
\end{proof}

\begin{lemma} 
$\mathcal{I},\sigma \vDash \forall x\, B
\quad\text{iff}\quad
\text{for all } d \in D,\;
\mathcal{I},\, \sigma[x \leftarrow d] \vDash B. $
\end{lemma}

\begin{lemma}
$\mathcal{I},\sigma \vDash \exists x\, B
\quad\text{iff}\quad
\text{for some } d \in D,\;
\mathcal{I},\, \sigma[x \leftarrow d] \vDash B.$
\end{lemma}
\bigskip
\textbf{Claim: } 
$\forall x\, (A(x)\land B(x)) \;\equiv\; (\forall x\,A(x)) \land (\forall x\,B(x)).$


\bigskip
\textbf{Proof.}

We show both directions of logical implication using the semantic definition of
\(\vDash\).

\paragraph{(1) $\forall x(A(x)\land B(x)) \;\Rightarrow\; (\forall xA(x))\land(\forall xB(x))$}

Let $\mathcal{I}$ be any interpretation with domain $D\neq\varnothing$, and let
$\sigma$ be any variable assignment.  
Assume

\[
\mathcal{I},\sigma \vDash \forall x\,(A(x)\land B(x)).
\]

By the semantics of the universal quantifier, this means:

\[
\text{for all } d\in D,\quad 
\mathcal{I},\sigma[x\mapsto d] \vDash A(x)\land B(x).
\]

By the semantics of conjunction:

\[
\text{for all } d\in D,\quad 
\mathcal{I},\sigma[x\mapsto d] \vDash A(x)
\quad\text{and}\quad
\mathcal{I},\sigma[x\mapsto d] \vDash B(x).
\]

Hence:

\[
\mathcal{I},\sigma \vDash \forall x\,A(x)
\qquad\text{and}\qquad
\mathcal{I},\sigma \vDash \forall x\,B(x).
\]

Thus, by semantics of $\land$,

\[
\mathcal{I},\sigma \vDash (\forall x\,A(x)) \land (\forall x\,B(x)).
\]

Since $\mathcal{I}$ and $\sigma$ were arbitrary, we obtain:

\[
\models\; \forall x(A(x)\land B(x)) \to 
\bigl((\forall x\,A(x))\land(\forall x\,B(x))\bigr).
\]


\paragraph{(2) $(\forall xA(x))\land(\forall xB(x)) \;\Rightarrow\; \forall x(A(x)\land B(x))$}

Again let $\mathcal{I}$ be any interpretation with domain $D\neq\varnothing$ and
$\sigma$ any assignment.  
Assume

\[
\mathcal{I},\sigma \vDash (\forall xA(x))\land(\forall xB(x)).
\]

By semantics of $\land$:

\[
\mathcal{I},\sigma \vDash \forall xA(x)
\qquad\text{and}\qquad
\mathcal{I},\sigma \vDash \forall xB(x).
\]

Thus, by semantics of $\forall$:

\[
\text{for all } d\in D,\quad 
\mathcal{I},\sigma[x\mapsto d] \vDash A(x)
\qquad\text{and}\qquad
\mathcal{I},\sigma[x\mapsto d] \vDash B(x).
\]

Therefore, for all $d\in D$,

\[
\mathcal{I},\sigma[x\mapsto d] \vDash A(x)\land B(x).
\]

Applying the semantics of $\forall$ again yields:

\[
\mathcal{I},\sigma \vDash \forall x\, (A(x)\land B(x)).
\]

Thus,

\[
\models\; \bigl((\forall xA(x))\land(\forall xB(x))\bigr)
     \to \forall x(A(x)\land B(x)).
\]


\textbf{Claim: } 
$\exists x\, (A(x)\lor B(x)) \;\equiv\; (\exists x\,A(x)) \lor (\exists x\,B(x)).$


\bigskip
\textbf{Proof.}

We prove both directions using the semantic definition of $\vDash$.

\paragraph{(1) $\exists x(A(x)\lor B(x)) \;\Rightarrow\; (\exists xA(x))\lor(\exists xB(x))$}

Let $\mathcal{I}$ be any interpretation with domain $D\neq\varnothing$, and
let $\sigma$ be any variable assignment.  
Assume

\[
\mathcal{I},\sigma \vDash \exists x\,(A(x)\lor B(x)).
\]

By the semantics of the existential quantifier, there exists some
$d\in D$ such that

\[
\mathcal{I},\sigma[x\mapsto d] \vDash A(x)\lor B(x).
\]

By the semantics of disjunction, this means:

\[
\mathcal{I},\sigma[x\mapsto d] \vDash A(x)
\quad\text{or}\quad
\mathcal{I},\sigma[x\mapsto d] \vDash B(x).
\]

\begin{itemize}
    \item If $\mathcal{I},\sigma[x\mapsto d] \vDash A(x)$, then by the
    semantics of $\exists$ we have
    $\mathcal{I},\sigma \vDash \exists x\,A(x)$.

    \item If $\mathcal{I},\sigma[x\mapsto d] \vDash B(x)$, then
    $\mathcal{I},\sigma \vDash \exists x\,B(x)$.
\end{itemize}

In either case, 

\[
\mathcal{I},\sigma \vDash (\exists x\,A(x)) \lor (\exists x\,B(x)).
\]

Since $\mathcal{I}$ and $\sigma$ were arbitrary, we obtain

\[
\models\; \exists x(A(x)\lor B(x)) 
\to \bigl((\exists x\,A(x))\lor(\exists x\,B(x))\bigr).
\]

\paragraph{(2) $(\exists xA(x))\lor(\exists xB(x)) \;\Rightarrow\; \exists x(A(x)\lor B(x))$}

Again let $\mathcal{I}$ be any interpretation with domain $D\neq\varnothing$
and $\sigma$ any assignment.  
Assume

\[
\mathcal{I},\sigma \vDash (\exists xA(x))\lor(\exists xB(x)).
\]

By the semantics of $\lor$,

\[
\mathcal{I},\sigma \vDash \exists xA(x)
\quad\text{or}\quad
\mathcal{I},\sigma \vDash \exists xB(x).
\]

\begin{itemize}
    \item Suppose $\mathcal{I},\sigma \vDash \exists xA(x)$.  
    Then there exists $d\in D$ such that
    $\mathcal{I},\sigma[x\mapsto d] \vDash A(x)$.
    Hence
    \[
       \mathcal{I},\sigma[x\mapsto d] \vDash A(x)\lor B(x),
    \]
    so by the semantics of $\exists$,
    \[
       \mathcal{I},\sigma \vDash \exists x(A(x)\lor B(x)).
    \]

    \item Suppose instead $\mathcal{I},\sigma \vDash \exists xB(x)$.  
    Then there exists $d\in D$ such that
    $\mathcal{I},\sigma[x\mapsto d] \vDash B(x)$, and hence again
    \[
       \mathcal{I},\sigma[x\mapsto d] \vDash A(x)\lor B(x),
    \]
    yielding
    \[
       \mathcal{I},\sigma \vDash \exists x(A(x)\lor B(x)).
    \]
\end{itemize}

Thus in all cases,

\[
\mathcal{I},\sigma \vDash \exists x(A(x)\lor B(x)).
\]

Therefore,

\[
\models\; 
\bigl((\exists xA(x))\lor(\exists xB(x))\bigr)
\to \exists x(A(x)\lor B(x)).
\]
\newpage
\begin{proposition}
In the structure $\mathcal I$ and assignment $\sigma$ (in class example),
\[
\mathcal I,\sigma \vDash \forall x\,\exists y\,\mathsf{lt}(x,y)
\quad\text{and}\quad
\mathcal I,\sigma \nvDash \exists y\,\forall x\,\mathsf{lt}(x,y).
\]
\end{proposition}

\begin{proof}
First formula. By the semantics of $\forall$,
\[
\mathcal I,\sigma \vDash \forall x\,\exists y\,\mathsf{lt}(x,y)
\]
iff for every $d\in D$,
\[
\mathcal I,\sigma[x\leftarrow d] \vDash \exists y\,\mathsf{lt}(x,y).
\]
Fix $d\in D$. By the semantics of $\exists$ this is equivalent to the existence
of some $e\in D$ such that
\[
\mathcal I,\sigma[x\leftarrow d][y\leftarrow e] \vDash \mathsf{lt}(x,y).
\]
Choose $e=d$. Then $d\le d$ in $\mathbb N$, so
$\mathsf{lt}^{\mathcal I}(d,d)$ holds, hence
\(\mathcal I,\sigma[x\leftarrow d][y\leftarrow d]\vDash \mathsf{lt}(x,y)\).
Thus for every $d$ such an $e$ exists, and therefore
\(\mathcal I,\sigma \vDash \forall x\,\exists y\,\mathsf{lt}(x,y)\).

\medskip
Second formula. Suppose, for contradiction, that
\[
\mathcal I,\sigma \vDash \exists y\,\forall x\,\mathsf{lt}(x,y).
\]
By the semantics of $\exists$, there is some $d\in D$ such that
\[
\mathcal I,\sigma[y\leftarrow d] \vDash \forall x\,\mathsf{lt}(x,y).
\]
By the semantics of $\forall$, this means that for every $e\in D$,
\[
\mathcal I,\sigma[y\leftarrow d][x\leftarrow e] \vDash \mathsf{lt}(x,y),
\]
i.e.\ $\mathsf{lt}^{\mathcal I}(e,d)$ holds for all $e\in D$.

Take $e=d+1\in\mathbb N$. Then $\mathsf{lt}^{\mathcal I}(d+1,d)$ would have
to hold, i.e.\ we would need $d+1\le d$, which is false in $\mathbb N$.
Hence
\[
\mathcal I,\sigma[y\leftarrow d][x\leftarrow d+1] \nvDash \mathsf{lt}(x,y),
\]
contradicting the assumption that
$\mathcal I,\sigma[y\leftarrow d] \vDash \forall x\,\mathsf{lt}(x,y)$.

Therefore no such $d$ exists, so
\[
\mathcal I,\sigma \nvDash \exists y\,\forall x\,\mathsf{lt}(x,y).
\]
\end{proof}

\begin{proposition}
Let $\Gamma$ be a set of first--order formulas and $A$ a formula.
Then
\[
\Gamma \vDash A \quad\Longleftrightarrow\quad
\Gamma \cup \{\neg A\} \text{ is unsatisfiable}.
\]
\end{proposition}

\begin{proof}
Recall:
\begin{itemize}
  \item $\mathcal I,\sigma \vDash \Gamma$ means $\mathcal I,\sigma \vDash B$ for every $B\in\Gamma$.
  \item $\Gamma \vDash A$ means: for all $\mathcal I,\sigma$, if $\mathcal I,\sigma \vDash \Gamma$
        then $\mathcal I,\sigma \vDash A$.
  \item $\Gamma$ is \emph{unsatisfiable} means: there are no $\mathcal I,\sigma$ with
        $\mathcal I,\sigma \vDash \Gamma$.
\end{itemize}

($\Rightarrow$) Assume $\Gamma \vDash A$. Suppose, for contradiction, that
$\Gamma\cup\{\neg A\}$ is satisfiable. Then there exist $\mathcal I,\sigma$ such that
$\mathcal I,\sigma \vDash \Gamma$ and $\mathcal I,\sigma \vDash \neg A$.
From $\Gamma \vDash A$ and $\mathcal I,\sigma \vDash \Gamma$ we obtain $\mathcal I,\sigma \vDash A$.
But $\mathcal I,\sigma \vDash \neg A$ iff $\mathcal I,\sigma \not\vDash A$, a contradiction.
Hence no such $\mathcal I,\sigma$ exist, so $\Gamma\cup\{\neg A\}$ is unsatisfiable.

\medskip
($\Leftarrow$) Assume $\Gamma\cup\{\neg A\}$ is unsatisfiable. To show $\Gamma \vDash A$,
let $\mathcal I,\sigma$ be arbitrary such that $\mathcal I,\sigma \vDash \Gamma$.
If $\mathcal I,\sigma \not\vDash A$, then by the negation clause we would have
$\mathcal I,\sigma \vDash \neg A$, and hence
\[
\mathcal I,\sigma \vDash \Gamma \cup \{\neg A\},
\]
contradicting the unsatisfiability of $\Gamma\cup\{\neg A\}$. Therefore
$\mathcal I,\sigma \vDash A$. Since $\mathcal I,\sigma$ was arbitrary, we conclude
$\Gamma \vDash A$.

Thus $\Gamma \vDash A$ iff $\Gamma\cup\{\neg A\}$ is unsatisfiable.
\end{proof}

Let
\[
A := \exists x \,\forall y\, p(x,f(y)),
\qquad
B := \forall y\, p(c,f(y)),
\]
where $c$ is a fresh constant symbol not appearing in $A$.
We prove that $A$ and $B$ are \emph{equi-satisfiable}.

\begin{theorem}
$A$ is satisfiable iff $B$ is satisfiable.
\end{theorem}

\begin{proof}
\textbf{($A$ satisfiable $\Rightarrow$ $B$ satisfiable).}
Assume $A$ is satisfiable.
Thus there exists a structure $\mathcal{I}$ and an assignment $\sigma$
such that
\[
\mathcal{I},\sigma \models \exists x\,\forall y\,p(x,f(y)).
\]
By the semantics of $\exists$, there exists an element $d \in D^{\mathcal{I}}$
such that
\[
\mathcal{I},\sigma[x := d] \models \forall y\,p(x,f(y)).
\]
Thus, for every $e \in D^{\mathcal{I}}$,
\[
\mathcal{I},\sigma[x:=d,y:=e] \models p(x,f(y)).
\]

Now define a new structure $\mathcal{I}'$ obtained from $\mathcal{I}$
by interpreting the new constant $c$ as $d$,
i.e.\ $c^{\mathcal{I}'} = d$, and leaving all other interpretations unchanged.
Then for every $e \in D^{\mathcal{I}'}$ we have
\[
\mathcal{I}',\sigma[y:=e] \models p(c,f(y)).
\]
Hence
\[
\mathcal{I}' \models \forall y\,p(c,f(y)),
\]
so $\mathcal{I}' \models B$, and therefore $B$ is satisfiable.

\bigskip

\textbf{($B$ satisfiable $\Rightarrow$ $A$ satisfiable).}
Assume that $B$ is satisfiable.
Thus there exists a structure $\mathcal{J}$ such that
\[
\mathcal{J} \models \forall y\, p(c,f(y)).
\]
Let $d := c^{\mathcal{J}}$ be the interpretation of $c$ in $\mathcal{J}$.
Then for every $e \in D^{\mathcal{J}}$,
\[
\mathcal{J},\sigma[y:=e] \models p(c,f(y)).
\]
Since $c^{\mathcal{J}} = d$, this means
\[
\mathcal{J},\sigma[x:=d,y:=e] \models p(x,f(y)).
\]

Therefore,
\[
\mathcal{J},\sigma[x:=d] \models \forall y\,p(x,f(y))
\]
for every assignment $\sigma$.
By choosing the witness $x := d$ we obtain
\[
\mathcal{J},\sigma \models \exists x\,\forall y\, p(x,f(y)).
\]
Thus $\mathcal{J},\sigma \models A$, and hence $A$ is satisfiable.

\bigskip

Since each direction holds, $A$ and $B$ are equi-satisfiable.
\end{proof}
\newpage
Let
\[
C := \forall y\,\exists x\,p(x,y),
\qquad
D := \forall y\,p(g(y),y),
\]
where $g$ is a fresh unary function symbol not appearing in $C$.
We prove that $C$ and $D$ are \emph{equi-satisfiable}.

\begin{theorem}
$C$ is satisfiable iff $D$ is satisfiable.
\end{theorem}

\begin{proof}
\textbf{($C$ satisfiable $\Rightarrow$ $D$ satisfiable).}
Assume that $C$ is satisfiable.
Then there exists a structure $\mathcal{I}$ and an assignment $\sigma$
such that
\[
\mathcal{I},\sigma \models \forall y\,\exists x\,p(x,y).
\]
By the semantics of $\forall$, this means that for every $d \in D^{\mathcal{I}}$,
\[
\mathcal{I},\sigma[y := d] \models \exists x\,p(x,y).
\]
By the semantics of $\exists$, for each $d \in D^{\mathcal{I}}$ there exists
some element $e \in D^{\mathcal{I}}$ (which may depend on $d$) such that
\[
\mathcal{I},\sigma[y := d, x := e] \models p(x,y).
\]

Define a function $h : D^{\mathcal{I}} \to D^{\mathcal{I}}$ by choosing, for
each $d \in D^{\mathcal{I}}$, one such witness $e$ and setting $h(d) := e$.
Now define a new structure $\mathcal{I}'$ that extends $\mathcal{I}$ by
interpreting the new function symbol $g$ as $h$, i.e.\ $g^{\mathcal{I}'} = h$,
and otherwise agrees with $\mathcal{I}$ on all symbols of the original language.

We claim that $\mathcal{I}' \models D$.
Let $d \in D^{\mathcal{I}'}$ be arbitrary.
By definition of $h$ and $g^{\mathcal{I}'}$, we have
$g^{\mathcal{I}'}(d) = h(d)$, and $h(d)$ is chosen so that
\[
\mathcal{I},\sigma[y := d, x := h(d)] \models p(x,y).
\]
Since $\mathcal{I}'$ and $\mathcal{I}$ agree on the interpretation of $p$,
it follows that
\[
\mathcal{I}',\sigma[y := d] \models p(g(y),y).
\]
Because $d$ was arbitrary, we conclude that
\[
\mathcal{I}' \models \forall y\,p(g(y),y),
\]
so $\mathcal{I}' \models D$, and therefore $D$ is satisfiable.

\bigskip

\textbf{($D$ satisfiable $\Rightarrow$ $C$ satisfiable).}
Conversely, suppose that $D$ is satisfiable.
Then there exists a structure $\mathcal{J}$ such that
\[
\mathcal{J} \models \forall y\,p(g(y),y).
\]
Let $h := g^{\mathcal{J}} : D^{\mathcal{J}} \to D^{\mathcal{J}}$ be the
interpretation of $g$ in $\mathcal{J}$.

We show that $\mathcal{J}$ is also a model of $C$ when viewed as a structure
for the original language (i.e.\ ignoring the extra symbol $g$ in the
statement of $C$).
Let $d \in D^{\mathcal{J}}$ be arbitrary.
Since
\[
\mathcal{J} \models \forall y\,p(g(y),y),
\]
we have, for every assignment $\sigma$,
\[
\mathcal{J},\sigma[y := d] \models p(g(y),y).
\]
Unfolding the term $g(y)$, this is the same as
\[
\mathcal{J},\sigma[y := d] \models p(h(d),y).
\]
Equivalently,
\[
\mathcal{J},\sigma[y := d, x := h(d)] \models p(x,y).
\]
Thus, for each $d \in D^{\mathcal{J}}$, there exists an element
$e = h(d) \in D^{\mathcal{J}}$ such that
\[
\mathcal{J},\sigma[y := d, x := e] \models p(x,y).
\]
By the semantics of $\exists$, this means
\[
\mathcal{J},\sigma[y := d] \models \exists x\,p(x,y)
\]
for every $d \in D^{\mathcal{J}}$.
By the semantics of $\forall$, we conclude that
\[
\mathcal{J},\sigma \models \forall y\,\exists x\,p(x,y),
\]
i.e.\ $\mathcal{J},\sigma \models C$.
Hence $C$ is satisfiable.

\bigskip

Since each direction holds, we conclude that $C$ and $D$ are equi-satisfiable.
\end{proof}

\end{document}
